\documentclass[11pt,a4paper]{article}

% Packages
\usepackage[left=1.5cm,right=1.5cm,top=2cm,bottom=2cm]{geometry}
\usepackage{fontspec}
\usepackage{xcolor}
\usepackage{hyperref}
\usepackage{titlesec}
\usepackage{enumitem}
\usepackage{graphicx}
\usepackage{fancyhdr}
\usepackage{tabularx}
\usepackage{newunicodechar}

% Farben
\definecolor{darkblue}{RGB}{0,47,108}
\definecolor{lightgray}{RGB}{100,100,100}

% Font
\setmainfont{Helvetica Neue}
\setsansfont{Helvetica Neue}
\renewcommand{\familydefault}{\sfdefault}

% Unicode-Pfeil definieren
\newunicodechar{→}{\ensuremath{\rightarrow}}

% Hyperref Setup
\hypersetup{
    colorlinks=true,
    linkcolor=darkblue,
    urlcolor=darkblue,
    pdfauthor={Keanu Fuchs},
    pdftitle={Lebenslauf Keanu Fuchs}
}

% Section Formatting
\titleformat{\section}
    {\Large\bfseries\color{darkblue}}
    {}
    {0em}
    {}
    [\titlerule]

\titlespacing*{\section}{0pt}{12pt}{6pt}

% Custom Commands
\newcommand{\cventry}[5]{
    \noindent
    \begin{minipage}[t]{0.75\textwidth}
        \textbf{#1}\\
        \textit{#2}
    \end{minipage}
    \hfill
    \begin{minipage}[t]{0.23\textwidth}
        \raggedleft\textcolor{lightgray}{#3}
    \end{minipage}
    \par\vspace{2pt}
    \noindent\textcolor{lightgray}{#4}
    \par\vspace{4pt}
    #5
    \par\vspace{8pt}
}

\newcommand{\cveducation}[4]{
    \noindent
    \begin{minipage}[t]{0.75\textwidth}
        \textbf{#1}\\
        \textit{#2}
    \end{minipage}
    \hfill
    \begin{minipage}[t]{0.23\textwidth}
        \raggedleft\textcolor{lightgray}{#3}
    \end{minipage}
    \par\vspace{2pt}
    #4
    \par\vspace{8pt}
}

% Header and Footer
\pagestyle{fancy}
\fancyhf{}
\renewcommand{\headrulewidth}{0pt}
\fancyfoot[L]{\small\textcolor{lightgray}{18. Januar 2026}}
\fancyfoot[C]{\small\textcolor{lightgray}{Keanu Fuchs}}
\fancyfoot[R]{\small\textcolor{lightgray}{\thepage}}

% List Settings
\setlist[itemize]{leftmargin=*,topsep=0pt,itemsep=2pt,parsep=0pt}

% Document
\begin{document}

% Disable page numbering for first page header
\thispagestyle{empty}

% Header
\noindent
\begin{minipage}[t]{0.65\textwidth}
    {\Huge\bfseries\color{darkblue}Keanu Fuchs}\\[4pt]
    {\Large\color{lightgray}System Engineer | Network Automation}\\[10pt]
    \textcolor{lightgray}{
        Schlesienstraße 5, 53119 Bonn\\
        \href{tel:+491746352391}{+49 174 6352391} | \href{mailto:kn.fchs@gmail.com}{kn.fchs@gmail.com}\\
        \href{https://linkedin.com/in/keanufuchs}{linkedin.com/in/keanufuchs} \\
        \href{https://github.com/keanufuchs}{github.com/keanufuchs} | \href{https://blog.spelk.de}{blog.spelk.de (Tech-Blog)}
    }
\end{minipage}
\hfill

\vspace{20pt}

% Enable fancy page style for subsequent pages
\pagestyle{fancy}

% Profil
\section*{Profil}
System Engineer und Dualer Student mit Fokus auf Netzwerkautomatisierung (Cisco) und Softwareentwicklung. Spezialisiert auf skalierbare Provisionierung (Day-0/Day-N) mittels Python, sowie die Integration moderner KI-Ansätze (LLM/RAG) in Network-Operations. Nachgewiesene Erfahrung in der technischen Leitung von Enterprise-Rollouts sowie der Entwicklung eigener Automatisierungstools im CI/CD-Kontext.
\vspace{10pt}

% Kompetenzen
\section*{Kompetenzen}
\begin{tabularx}{\textwidth}{@{}X@{\hspace{2em}}X@{}}
    \textbf{Network Automation} & \textbf{APIs \& Protokolle}\\
    Cisco Catalyst Center, Plug \& Play, Day-0/Day-N, Templates (Jinja2), InfraHub (SSoT) & REST APIs, NETCONF/YANG, Cisco APIs (Catalyst Center, FMC, Meraki)\\[8pt]
    
    \textbf{Cisco Networking} & \textbf{Automation \& Testing}\\
    Catalyst Switching, Enterprise Rollouts, Troubleshooting, Meraki, Firepower & Python, PyATS, Ansible, Paramiko, Flask, Pydantic\\[8pt]
    
    \textbf{DevOps \& Development} & \textbf{Data Handling}\\
    Git, GitLab CI/CD, Docker, GitLab Runner, Vue.js, Nuxt & Excel-Validierung, YAML/JSON Transformation, Terraform Integration\\
\end{tabularx}

\vspace{10pt}

% Berufserfahrung
\section*{Berufserfahrung}

\cventry{Dualer Student Informatik (B.Sc.) \& System Engineer}{Bechtle IT-Systemhaus Bonn/Köln, Bonn}{Sep. 2023 – heute}{}{
\begin{itemize}
    \item Netzwerkautomatisierung im Cisco-Enterprise-Umfeld (Catalyst Center, Python, REST, NETCONF/YANG)
    \item Automatisierung von Provisionierung \& Lifecycle-Prozessen (Day-0/Day-N)
    \item API-basierte Integrationen \& Tools im Bereich Switching und Security
    \item CI/CD-basierte Umsetzung über GitLab Pipelines inkl. Docker-Containerisierung
\end{itemize}
\textbf{Projekt-Highlights (anonymisiert):}
\begin{itemize}
    \item Enterprise Rollout (2023–2024): 11 Logistikzentren + 8 Büro-/Kleinstandorte, 10–30 Devices/Standort, technischer Lead
    \item Industrie-Rollouts (2024–heute): 5 Standorte, Blueprints/Standards, eigenverantwortlich
    \item Catalyst Center API Extensions (Feb.–Sep. 2024): Python REST API
    \item Automatisierte Datenmigration \& Validierung (Excel zu YAML) (Sep.–Okt. 2024): Python/Pandas
    \item Firewall Automation Cisco FMC API (Jan. 2025): Python REST API
    \item Cisco Catalyst Center PnP-Emulation Tool (Mär.–Jun. 2025): Vue.js/Nuxt, Python Flask, Ansible
    \item Confluence → Pydantic → Terraform Pipeline (Jul.–Okt. 2025): Python/Pydantic
    \item Hardening \& Compliance Automation (Nov. 2025–heute): NETCONF/YANG, PyATS
\end{itemize}
}

\cventry{Junior System Engineer}{Bechtle IT-Systemhaus Bonn/Köln, Bonn}{Jun. 2023 – Sep. 2023}{}{
\begin{itemize}
    \item Verantwortung im Cisco Catalyst Center für Provisionierung \& Lifecycle-Prozesse
    \item Plug-and-Play Rollouts sowie Day-0/Day-N Automatisierung (Templates/Jinja2)
    \item Unterstützung bei Standardisierung, Umsetzung und Troubleshooting
\end{itemize}
}

\cventry{Fachinformatiker für Systemintegration (Ausbildung)}{Bechtle IT-Systemhaus Bonn/Köln, Bonn}{Sep. 2021 – Jun. 2023}{}{
\begin{itemize}
    \item Schwerpunkt Cisco Networking inkl. Lifecycle-Prozesse
    \item Mitarbeit in Enterprise-Netzwerkumgebungen
    \item \textbf{Abschlussprojekt:} Automatisierter Geräteaustausch \& Provisionierung über Cisco Catalyst Center
\end{itemize}
}

% Ausbildung
\section*{Ausbildung}

\cveducation{Bachelor of Science (B.Sc.) Informatik (dual)}{Rheinische Hochschule Köln}{Sep. 2023 – Jul. 2026}{
\begin{itemize}
    \item Schwerpunkt: AI Solution Development und Network Automation
    
    % ÄNDERUNG: Bachelorarbeit visuell hervorgehoben durch \vspace und Fettdruck
    \item[] \vspace{2pt} \textbf{Bachelorarbeit (in Arbeit):} \\
    \textit{„Von natürlicher Sprache zu Netzwerk-Konfigurationen: Ein LLM-gestütztes Intent-Based Networking-System mit MCP-basierter Single Source of Truth“}
    
    % Kurze Erklärung der Tech-Stack macht es greifbarer:
    \item[] \footnotesize{\textcolor{darkblue}{\textbf{Tech Stack:}} Python, LangChain, InfraHub (SSoT), Ollama/OpenAI API. Entwicklung eines Agenten, der Netzwerkinfrastruktur basierend auf natürlicher Sprache provisioniert.}
\end{itemize}
}

\cveducation{Fachinformatiker für Systemintegration (IHK)}{Heinrich-Hertz-Europakolleg Bonn}{Sep. 2021 – Jun. 2023}{
\begin{itemize}
    \item Ausbildung mit Schwerpunkt Netzwerktechnik und Cisco-Systeme
\end{itemize}
}

\cveducation{Allgemeine Hochschulreife (Abitur)}{Städtisches Gymnasium Rheinbach}{Aug. 2013 – Jun. 2021}{
}

% Zertifikate
\section*{Zertifikate}
\noindent
\textbf{Cisco CCNA (Cisco Certified Network Associate)} \hfill \textcolor{lightgray}{Feb. 2024 – Feb. 2027}

\vspace{10pt}

% Sprachen
\section*{Sprachen}
\begin{tabularx}{\textwidth}{@{}X@{\hspace{2em}}X@{}}
    \textbf{Deutsch:} Muttersprache & \textbf{Englisch:} Sehr gut (Wort und Schrift)\\[4pt]
    \textbf{Französisch:} Grundkenntnisse & \\
\end{tabularx}

\vspace{10pt}

% Interessen & Private Projekte
\section*{Interessen \& Private Projekte}
\textbf{Tech-Blog \& HomeLab (blog.spelk.de):} Dokumentation technischer Lösungen (z.B. Cisco Catalyst Checklisten, Hardware-Recovery) sowie Server-Architekturing (On-Demand Proxy-Scaling mit Velocity/Limbo).\\
\textbf{Sonstiges:} Fitness (seit 2019), Fußball (FC Hertha Bonn, 5 Jahre), Waldhorn (Musikschule Bonn, 3 Jahre)

\end{document}
